\documentclass{article}

% NeurIPS 2025 style. Use \usepackage[preprint]{neurips_2025} for a preprint,
% or \usepackage[final]{neurips_2025} for a camera-ready version.
% The only supported style file for NeurIPS 2026 is neurips_2026.sty,
% obtainable from neurips.cc. Options: final, preprint, nonatbib.
% Use [preprint] for arXiv, no option for anonymous submission.
\usepackage{neurips_2026}

\usepackage[utf8]{inputenc}
\usepackage[T1]{fontenc}
\usepackage{hyperref}
\usepackage{url}
\usepackage{booktabs}   % required by the instructions for table rules
\usepackage{amsfonts}
\usepackage{amsmath}
\usepackage{nicefrac}
\usepackage{microtype}
\usepackage{graphicx}
\usepackage{xcolor}   % \answerYes/No/NA use \textcolor

\title{You Can't Pretrain Your Way Out of a Missing Negative Class:\\
Nurse Stress Detection from Wearables}

\author{%
  Anonymous Author(s) \\
  Affiliation \\
  Address \\
  \texttt{email} \\
}

\begin{document}

\maketitle

\begin{abstract}
We set out to apply self-supervised pretraining to nurse stress detection,
on a dataset where 90.9\% of recorded sensor time carries no label
at all. The appeal was obvious: two hundred times more unlabelled than labelled
data, and a natural pretext task. Establishing the prerequisites consumed the
study, and then pretraining did not help. Pretraining requires a negative
class, and no participant ever recorded being unstressed: all 358 survey
entries mark an episode a nurse flagged, and the rating is its severity. We show that the labels do
mark a physiological state, but in one channel only: skin conductance rises at
reported onset in 64.3\% of episodes ($p = 0.0009$), while heart
rate is indistinguishable from chance ($p = 0.80$). We construct a negative
class by activity-matched sampling and verify by falsification that the
resulting classifier detects neither movement nor shift schedule (both
AUC $0.51$). The baseline reaches AUC $0.765$, far above a within-subject
permutation null ($0.497 \pm 0.010$), yet recovers only 10.6\% of
episodes at one false alarm per worn hour. A linear model reaches within
$0.015$ AUC of a boosted one, so capacity is not the constraint; three
treatments of the labels, including non-negative positive-unlabelled risk
estimation, all fail to improve on it. Masked-reconstruction pretraining over
1,238 hours then fails as well: a pretrained encoder matches a randomly
initialised one at every label count from 25 to 138 episodes, and both trail
the hand-crafted baseline by $0.11$ AUC. Severity prediction, which needs no
constructed negatives at all, is also at chance. Two findings run against
expectation: careful negative construction improves rather than inflates
performance, and a published accuracy figure on this dataset is difficult to
reconcile with subject-grouped evaluation.
\end{abstract}

\section{Introduction}

Occupational stress among nurses is associated with burnout, turnover, and
adverse patient events. Wrist-worn sensors offer a continuous alternative to
retrospective self-report, and a public dataset exists for studying it:
Hosseini et al.\ \cite{hosseini2022} monitored 15 nurses for one week each
during the COVID-19 outbreak, recording electrodermal activity, heart rate,
skin temperature, and acceleration, with periodic smartphone-administered
stress reports.

The dataset has a property that makes self-supervision attractive. Only
9.1\% of the 1251.7 hours of recording falls inside a rated
episode. The remaining 90.9\% is unlabelled physiology, which is
exactly the regime where masked-reconstruction pretraining is supposed to help:
learn what ordinary physiology looks like from everything, then fine-tune on
the little that is labelled. That was our plan.

Two prerequisites stopped it. Self-supervision does not manufacture a negative
class, and this dataset has none: every one of the 358 survey rows marks an
episode a participant flagged, and the number attached is its severity, not its
presence. Nor does self-supervision have anything to improve on until a
baseline exists. Building both turned out to be the study.

We therefore report the groundwork rather than the method we intended, and we
organise it as a sequence of eliminations against one question:

\begin{quote}
\emph{What would have to be true for representation learning to help here?}
\end{quote}

Section \ref{sec:labels} asks whether the labels mark anything physiological.
Section \ref{sec:negatives} asks whether the comparison group we build leaks a
confound. Section \ref{sec:baseline} asks how far standard features and a
standard learner get, and whether the pipeline manufactures its own signal.
Section \ref{sec:pu} asks whether a different treatment of the labels helps.
Each answer narrows what remains.

\paragraph{Contributions.}
\begin{enumerate}
\item An event-triggered analysis showing that only one of four channels
      responds at reported onset, which reorders the feature priorities for
      this dataset.
\item A negative-class construction with falsification probes, and the finding
      that careful construction \emph{improves} performance rather than
      inflating it, contrary to the usual warning.
\item A decision rule fixed before results were seen, governing whether to
      escalate to representation learning, together with the measurement that
      answered it and an explicit statement of what that measurement does not
      settle.
\item Two controls, a second learner and a within-subject permutation null,
      that separate what the data supports from what the pipeline produces.
\end{enumerate}

\section{Data and units of analysis}
\label{sec:data}

\subsection{Source}

We use the original Dryad deposit rather than the merged convenience files in
common circulation. The merged version resamples every channel onto a shared
grid, which discards the blood volume pulse and inter-beat interval streams and
removes the only route to heart rate variability. Channels are read at their
native rates: acceleration at 32 Hz, blood volume pulse at
64 Hz, electrodermal activity and temperature at 4 Hz, and
heart rate at 1 Hz, taken from the file headers. Published
descriptions giving 72 Hz and 10 Hz for pulse and
temperature disagree with the files themselves, and we follow the files.

The archive contains 609 sessions totalling 1251.7 hours across 250
person-days. The median session is 1.17 hours rather than a full shift,
so a participant's week is fragmented across roughly 40 short recordings.

\subsection{Clock alignment}

Sensor timestamps are UTC; survey timestamps are bare local wall-clock with no
zone recorded. Choosing wrongly places nearly every episode outside its own
recording and raises no error, so we resolved it empirically by counting how
many rated episodes start inside a session belonging to the same participant.
Treating them as UTC places 50 of 245 (20.4\%); as daylight-saving-aware local
time, 212 (86.5\%). Thirty-three rated episodes (13.5\%) fall outside every
session even so, and are set aside.

\subsection{Window, label, and the metric}
\label{sec:metric}

The unit of prediction is a 120 s window with a 60 s
hop. Two arguments fix this length. No labelled episode is shorter than
120 s, and a 300 s window would discard 15 episodes
outright. Independently, ultra-short heart rate variability work finds
120 s sufficient for time-domain measures \cite{munoz2015}, with
root mean square of successive differences validated from recordings as short
as 15 s \cite{orini2023}; the classical five-minute figure is a
recording convention rather than a requirement.

A window is \textbf{positive} if at least half of it falls inside a
high-severity episode. A window is \textbf{negative} only if it survives the
construction in Section \ref{sec:negatives}. Everything else is
\textbf{excluded}, not set to zero.

\paragraph{The primary metric, defined before use.} We report
\emph{episode-level recall at a fixed alarm rate}, and both halves need
stating.

An \textbf{episode is detected} if at least half of its constituent windows
fire. The fraction is fixed at one half in advance; we do not tune it.

The \textbf{alarm rate} is expressed per hour of worn time, not per hour of
sampled negative time. These differ substantially. Our scored sample is
43.6\% positive by construction, whereas true window prevalence
over all worn windows is 6.49\%, so counting alarms against the
sampled negatives would understate the deployed rate roughly sevenfold. We
therefore fix the operating point on the false positive rate, which is
prevalence-invariant, and translate: at 6.49\% prevalence there
are $28.1$ negative windows per worn hour, so one alarm per hour corresponds to
a false positive rate of $0.0356$. All operating points below are reported this
way.

\subsection{The structural problem}

The survey contains 358 rows, of which 245 carry a severity rating, split 46
low, 20 medium, 179 high. \textbf{None of them means ``not stressed.''} Every
row is an episode a participant actively flagged, and the rating is its
severity. The majority labelled class is therefore the most severe one, and
nine of fifteen participants have no medium-severity episodes at all.

Two consequences follow. Three-class modelling is unworkable, so we collapse to
high-severity against a constructed baseline. And the comparison group does not
exist in the data, making this a positive-unlabelled problem \cite{elkan2008}
rather than ordinary binary classification.

\subsection{Exclusions and cohort}

Exclusions are applied in a declared order, because they do not commute:
measuring session length before removing non-wear time retains 517 sessions,
and measuring it after retains 504. Thirteen sessions turn on rule order alone.

We remove short sessions, non-wear windows, sessions with a dead electrodermal
sensor, unrated episodes from both classes, duplicate survey rows, and low- and
medium-severity episodes; twelve overlapping episode pairs are merged to their
outer span.

We further remove five participants. One has no eligible comparison windows at
all. Four have median skin conductance between 0.07  and
0.11 $\mu$S, at or below the sensor floor. A low median alone
does not disqualify a participant, since a small but varying signal remains
informative after normalisation. The deciding statistic is relative variation,
the interquartile range divided by the median: these four score $0.44$ to
$0.72$, while the remaining eleven score $1.04$ to $5.14$, with no overlap.
Their signal is flat, not merely small.

The retained cohort is 10 participants, 24{,}334 windows, and 138
high-severity episodes with sensor coverage. Three counts describe this dataset
and only two are denominators: the window count is a compute statistic, the
honest denominator for detecting an episode is 138, and for generalising to a
new person it is 10.

\section{Do the labels mark anything physiological?}
\label{sec:labels}

Before building a classifier we ask whether there is anything to classify. We
align every high-severity episode at its reported onset, normalise each channel
against the participant's own preceding hour, and compare the first ten minutes
after onset against the preceding thirty.

\begin{table}[t]
\caption{Event-triggered response at reported onset, 140 episodes with usable
coverage. Robust statistics only; see the note below on why.}
\label{tab:onset}
\centering
\begin{tabular}{lrrrr}
\toprule
Channel & Median $\Delta z$ & Episodes positive & Sign test & Median raw shift \\
\midrule
Electrodermal activity & $+0.568$ & 90/140 (64.3\%) & $p = 0.0009$ &
  $+0.64\,\mu\mathrm{S}$ \\
Heart rate & $-0.067$ & 68/140 (48.6\%) & $p = 0.80$ & $+0.47$ bpm \\
\bottomrule
\end{tabular}
\end{table}

Table \ref{tab:onset} gives the result. Skin conductance rises at onset in
64.3\% of episodes, a shift that survives both a sign test and a
Wilcoxon signed-rank test ($p = 3.3 \times 10^{-6}$), and the raw median shift
of 0.64 $\mu$S sits on a baseline of 0.68 $\mu$S,
roughly a doubling. Heart rate rises in 48.6\% of episodes, which
is indistinguishable from chance.

\paragraph{Why we report medians.} The pooled mean tells a different and false
story. It reports an electrodermal shift of $+9.75$ normalised units, of which
63\% is contributed by five episodes out of 140. The cause is
numerical rather than physiological: the trailing interquartile range used as a
denominator approaches zero during flat stretches, and the largest single value
reached $+1600$ where a plausible score is near $\pm 3$. We bound the
denominator at its empirical first percentile, clip the output, and fix the use
of medians and sign tests in advance, because the same data yields opposite
headlines otherwise.

This analysis is descriptive and uses all episodes; it informs which channels
enter the feature set but is not itself a predictive claim, and no model is
selected on it.

The result reorders the rest of the study. The problem is well posed, and it is
well posed in one channel. Any subsequent result driven by heart rate deserves
suspicion until explained.

\section{Constructing a comparison group}
\label{sec:negatives}

\subsection{Eligibility}

A candidate negative window must satisfy four conditions: it lies in a session
of at least thirty minutes; it falls on a participant-day containing at least
one reported episode; it lies at least thirty minutes from the boundary of any
reported episode, including unrated ones; and it is not itself positive.

The second condition carries most of the argument. On a day when a participant
filed nothing, silence is uninformative: she may have been calm, or too busy to
open the application. On a day when she filed at least one report, she was
demonstrably engaging with the survey, so silence during the remainder of that
day is evidence rather than absence of evidence. This condition is also the
most expensive, reducing 24{,}334 windows to 16{,}214. All four together leave
7{,}955 candidate windows, or 265.2 hours.

\subsection{Matching rather than filtering}

Nurses move more when busy, and busy periods are stressful, so movement
correlates with stress without being a stress signal. Selecting quiet periods
as negatives would produce two classes differing in both stress and movement,
and a classifier would separate them by detecting movement.

We therefore match rather than filter. We bin the positive class into deciles
of mean acceleration magnitude, and for each bin draw negatives from the same
bin, \emph{within the same participant}. Global matching would draw negatives
disproportionately from high-coverage participants, making participant identity
a partial label proxy along exactly the axis the evaluation holds out.

\subsection{The negative ratio}

We fix the negative-to-positive ratio uniformly rather than letting each
participant supply its maximum, because class balance determines where the
decision threshold belongs. Table \ref{tab:f1} makes the point concretely: one
fixed model, recall $0.80$ and specificity $0.80$ throughout, scored on two
folds differing only in prevalence.

\begin{table}[t]
\caption{The same model on 100 windows at two prevalences. Recall and
specificity are held fixed; only the class balance changes. $F_1$ falls from
$0.727$ to $0.195$ because the false positives are drawn from a much larger
negative pool.}
\label{tab:f1}
\centering
\begin{tabular}{lrr}
\toprule
 & 2:1 fold (33.3\% positive) & 32:1 fold (3.0\% positive) \\
\midrule
Precision & 0.667 & 0.111 \\
Recall    & 0.800 & 0.800 \\
$F_1$     & \textbf{0.727} & \textbf{0.195} \\
\bottomrule
\end{tabular}
\end{table}

Exact uniformity is unattainable. The binding constraint is a single cell, one
participant's highest-activity bin containing 13 positives and 7 candidates, so
matching every participant exactly would require a ratio below one. We report a
sensitivity analysis across four arms in Table \ref{tab:sens} and fix
$\text{ratio} = 1.31$, which halves the residual spread from $1.64\times$ to
$1.28\times$. We note that the alternative arm at $2.14$ yields marginally
higher AUC; we prefer $1.31$ because comparable per-fold thresholds matter more
here than a difference of $0.002$, and we report both rather than choosing
silently.

Negatives are drawn under three random seeds. The seed-to-seed spread is the
noise floor: any later comparison smaller than it is not a comparison.

\subsection{Falsification}

We then build models we want to fail. Each is given only one suspicious signal,
using the same learner as the baseline so that a null result is attributable to
absence of signal rather than to a weaker model.

\begin{table}[t]
\caption{Falsification probes and ratio sensitivity, mean of three seeds. All
probes pass in all arms. Area under the ROC curve is prevalence-invariant, so
arms with different denominators remain comparable.}
\label{tab:sens}
\centering
\small
\begin{tabular}{lrrrrrr}
\toprule
Arm & Participants & Balance spread & ACC only & Time only & EDA only & Full \\
\midrule
All 10, ratio 1.31   & 10 & $1.28\times$ & 0.513 & 0.507 & 0.768 & 0.749 \\
All 10, ratio 2.14   & 10 & $1.64\times$ & 0.511 & 0.510 & 0.770 & 0.736 \\
Drop one, ratio 1.31 &  9 & $1.07\times$ & 0.521 & 0.533 & 0.762 & 0.739 \\
Drop one, ratio 2.25 &  9 & $1.30\times$ & 0.517 & 0.525 & 0.765 & 0.726 \\
\bottomrule
\end{tabular}
\end{table}

A model given only acceleration reaches AUC $0.513$, and one given only the
hour of day reaches $0.507$, against a chance value of $0.500$. Both hold in
all four arms. The constructed comparison group therefore does not leak
movement or shift schedule.

Participant identity is recoverable from the same features at $0.768$ accuracy
against a chance value of $0.100$. We record this as a measurement rather than
a failure, since physiology genuinely differs between people; it is the
reference against which any learned representation would have to be compared.

\section{Baseline and controls}
\label{sec:baseline}

Features are extracted at native rates and then aggregated to the window.
Detecting skin conductance responses on a 1 Hz downsample loses
roughly 40\% of them, so the common 1 Hz grid is used
only for the label join and coarse features. The feature set comprises tonic
level and slope, phasic variability, response count and mean amplitude for
electrodermal activity; mean, variability, slope, maximum and deviation from a
trailing median for heart rate; magnitude statistics and a movement fraction
for acceleration; slope and deviation from a trailing median for temperature,
never its absolute value, which has an intraclass correlation of $0.52$ with
participant identity and acts as a fingerprint; and hour of day with position
in session.

Heart rate variability is excluded from the baseline. Although the inter-beat
interval file is present in every session, only 5.9\% of
120 s windows contain a run of clean consecutive beats long enough
to compute it, and the archive holds only 39.5 hours of usable beat data
out of 1251.7 hours. Reported session-level availability of
48\% is eight times more flattering than the window-level figure
that governs a per-window feature.

Continuous channels are normalised against a trailing sixty-minute robust
baseline within session, which uses only past data and is therefore both
leakage-free and reproducible at deployment. The model is gradient boosting
with balanced class weights, evaluated leave-one-participant-out across three
negative-sampling seeds.

\begin{table}[t]
\caption{Baseline operating points. The alarm rate follows the translation in
Section \ref{sec:metric}. The seed range on episode recall at one alarm per
hour is comparable to the estimate itself, and we quote it wherever the
estimate appears.}
\label{tab:baseline}
\centering
\begin{tabular}{lrrrr}
\toprule
Alarms per worn hour & False positive rate & Window recall &
Episode recall & Seed range \\
\midrule
0.5 & 0.018 & 0.087 & 0.029 & 0.000 \\
1.0 & 0.036 & 0.155 & \textbf{0.106} & 0.072 \\
2.0 & 0.071 & 0.292 & 0.280 & 0.058 \\
\bottomrule
\end{tabular}
\end{table}

The model attains AUC $0.7654$ with a seed range of $0.0095$. Table
\ref{tab:baseline} translates this into operating points.

\paragraph{The gap is the finding.} An AUC of $0.765$ means the model reliably
ranks stressed windows above unstressed ones. At one false alarm per worn hour
it recovers 10.6\% of episodes, with a seed range of $0.072$ that
is comparable to the estimate itself. AUC integrates over false positive rates
including regions no deployment would tolerate; once constrained to the usable
corner, little remains. Two of ten participants detect zero episodes at any
usable threshold, and one is below chance at AUC $0.388$. We report per fold
and never average; fold AUC ranges from $0.388$ to $0.885$.

\subsection{Two controls}

\begin{table}[t]
\caption{Controls. The permutation null shuffles labels within participant,
preserving class balance, fold structure and feature distributions while
destroying only the label-feature correspondence.}
\label{tab:controls}
\centering
\begin{tabular}{lrr}
\toprule
 & AUC & Episode recall at 1 alarm/hour \\
\midrule
Gradient boosting        & 0.7655 & 0.1063 \\
Logistic regression      & 0.7502 & 0.0942 \\
\midrule
Permutation null (mean of 20) & 0.4975 ($\pm 0.0099$) & --- \\
Permutation null (maximum)    & 0.5186 & --- \\
\bottomrule
\end{tabular}
\end{table}

Table \ref{tab:controls} reports two controls that separate what the data
supports from what the pipeline produces.

Shuffling labels within participant yields AUC $0.4975 \pm 0.0099$ across
twenty permutations, with a maximum of $0.5186$. The observed value exceeds
every null draw, roughly twenty-six standard deviations above the null mean.
The pipeline does not manufacture its own signal.

A penalised logistic regression on identical folds and features reaches AUC
$0.7502$, within $0.0153$ of the boosted model. A linear model with almost no
capacity to overfit arrives at nearly the same place, so model capacity is not
the binding constraint.

\section{Does a different treatment of the labels help?}
\label{sec:pu}

We fixed a decision rule before computing the following, and record it here as
written:

\begin{quote}
Proceed to representation learning only if a different label treatment moves
episode recall at one alarm per worn hour by more than the baseline seed range
of $0.0724$.
\end{quote}

The threshold is the baseline's own measured seed range rather than a chosen
constant. We note that this makes the rule generous in one direction: a noisier
baseline sets a higher bar and is easier to leave uncleared.

We compare three treatments on identical folds, windows, and features: the
curated negatives of Section \ref{sec:negatives}; a naive comparator treating
all unlabelled time as negative, with no same-day rule, guard band, or
matching; and non-negative positive-unlabelled risk estimation
\cite{kiryo2017} fitted participant-stratified, with the class prior bracketed
rather than estimated, since it is not identifiable from positive-unlabelled
data.

\begin{table}[t]
\caption{Label treatments on identical folds. Both reported metrics are
prevalence-invariant, so the naive arm scored over 24{,}334 windows and the
curated arm over 3{,}622 remain comparable. The positive-unlabelled rows use a
fixed training budget without an inner validation split and are marked
accordingly; they are not tuned to the same standard as the baseline.}
\label{tab:pu}
\centering
\begin{tabular}{llrr}
\toprule
Treatment & Class prior & AUC & Episode recall at 1 alarm/hour \\
\midrule
Curated negatives (baseline) & --- & \textbf{0.7654} & \textbf{0.1063} \\
Naive, all unlabelled negative & --- & 0.6925 & 0.0652 \\
Non-negative PU (untuned) & 0.05 & 0.6616 & 0.0435 \\
Non-negative PU (untuned) & 0.10 & 0.6933 & 0.0652 \\
Non-negative PU (untuned) & 0.20 & 0.7132 & 0.0507 \\
Non-negative PU (untuned) & 0.30 & 0.7167 & 0.0580 \\
\bottomrule
\end{tabular}
\end{table}

The threshold to clear was $0.1787$. The best alternative reached $0.0652$, a
change of $-0.0411$. \textbf{The rule was not met, and every alternative is
worse than the baseline rather than merely insufficiently better.}

\paragraph{Curation improved rather than inflated.} This comparison is normally
run to detect inflation: naive negatives include off-shift and sleeping time,
the contrast becomes trivial, and the score rises. Here the curated
construction beats naive by $0.073$ AUC and $0.041$ episode recall. The likely
mechanism is the guard band, since naive negatives include windows adjacent to
episodes, which are physiologically similar to positives and act as label noise
during training. We have not isolated the guard band from the same-day rule,
so we offer this as the plausible mechanism rather than a demonstrated one.

\paragraph{The prior does not rescue it.} AUC rises monotonically with the
class prior, from $0.662$ to $0.717$, which is the shape one could present as
promising. Episode recall at the deployable operating point does not follow,
remaining between $0.044$ and $0.065$ throughout. We report the prior-indexed
results in full rather than the trend alone.

\section{Does a learned representation help?}
\label{sec:ssl}

The decision rule in Section \ref{sec:pu} tested how the labels are
\emph{used}. Self-supervision asks a different question: whether a different
\emph{representation} helps. We ran it, past our own stopping point, because
the two axes are not the same and a null on one does not settle the other.

\paragraph{Pretraining.} A dilated 1D convolutional encoder is trained by
masked reconstruction over all 609 sessions and all 15 participants, including
the five excluded from the supervised cohort, since the pretext task needs no
labels. The corpus is 1,238 usable hours at 1 Hz across five causally
normalised channels, yielding 68k windows. Mask spans of 30 to 120 s are fixed
a priori rather than tuned on downstream performance, which would leak through
the selection process. Contrastive learning was rejected deliberately: its
standard augmentations include amplitude scaling, and electrodermal amplitude
is the signal Section \ref{sec:labels} identified.

\paragraph{The identity gate.} A linear probe on frozen embeddings recovers
participant identity at $0.168$, against chance $0.067$ and the $0.768$
reachable from hand-crafted features. The encoder did not learn identity, so
the phase proceeds. The competing reading, that $0.168$ indicates a weak
representation rather than a clean one, is settled by the sweep below.

\begin{table}[t]
\caption{Label-efficiency sweep. Identical folds and windows to Table
\ref{tab:baseline}. Training positives are restricted to a sampled subset of
episodes; the test fold is never subsampled. All values are transductive upper
bounds, since the encoder saw every participant during pretraining.}
\label{tab:ssl}
\centering
\begin{tabular}{lrrrr}
\toprule
Episodes used & Linear probe & Full fine-tune & Random init & Probe $-$ random \\
\midrule
25  & 0.6157 & 0.6382 & 0.6398 & $-0.024$ \\
50  & 0.6342 & 0.6438 & 0.6367 & $-0.003$ \\
100 & 0.6486 & 0.6602 & 0.6604 & $-0.012$ \\
138 & 0.6467 & 0.6570 & 0.6495 & $-0.003$ \\
\midrule
\multicolumn{4}{l}{\emph{Hand-crafted baseline, Section \ref{sec:baseline}}}
  & \textbf{0.7654} \\
\bottomrule
\end{tabular}
\end{table}

\paragraph{No label efficiency, and no benefit.} Table \ref{tab:ssl} gives
two negatives. Pretrained minus randomly initialised is between $-0.024$ and
$-0.003$ at every label count, so there is no separation at low counts, which
was the entire hypothesis. And every arm trails the hand-crafted baseline by
roughly $0.11$ AUC.

That the randomly initialised arm matches the pretrained one is the
informative part. The architecture is not the problem, reaching $0.65$ from
random weights; the pretext task simply added nothing to it. This also
resolves the ambiguity in the identity probe: the representation is weak
rather than clean. And because the encoder saw every participant, these are
upper bounds, so the honest comparison is less favourable still.

\paragraph{Severity, where no negatives are constructed.} Detection could in
principle be blamed on our constructed comparison group. Severity cannot: it is
a three-level ordinal task over the 186 rated episodes with sensor coverage,
and it needs no negatives at all. We fit cumulative binary links rather than a
three-class softmax, because 14 medium-severity episodes across the cohort and
six of ten participants with none make the middle class unlearnable, and
because the rating scales are not comparable between people.

Both links sit at chance: $0.4964$ for severity $\geq 1$ and $0.5417$ for
severity $\geq 2$, with per-participant values from $0.300$ to $0.808$. Nor is
severity a restatement of episode length, since the point-biserial correlation
with duration is $-0.03$ and $-0.06$, neither significant. The one task in this
study whose numbers are not conditional on a construction fails too, which
removes the negative construction as the explanation for the rest.

\section{Discussion}

\subsection{What the eliminations leave, and what they do not}

The labels mark a real physiological state, in one channel. The comparison
group leaks neither movement nor schedule. The pipeline does not manufacture
signal, since a within-participant permutation null sits at chance. Model
capacity is not the constraint, since a linear model matches a boosted one.
Three treatments of the labels do not improve on the baseline, and two make it
worse.

Section \ref{sec:ssl} closes the remaining axis. A learned representation does
not move the result either, and neither does severity prediction, which needs
no constructed negatives at all.

What remains is the data: 138 episodes from 10 participants; labels that are
retrospective self-reports of intervals, treated here as uniform states
although a thirty-minute report is unlikely to describe thirty uniform
minutes; a constructed rather than observed negative class; and one responsive
channel out of four.

\subsection{On the reported accuracy in prior work}

Mathur et al.\ \cite{mathur2024} report weighted $F_1 = 0.99$ on this dataset.
We do not reproduce this and think it is difficult to reconcile with
participant-grouped evaluation. Adjacent windows in these recordings share raw
samples and participant-specific baselines; a split that does not respect
participant boundaries places near-duplicates on both sides. We are not in a
position to diagnose another group's protocol from a published summary, and we
note the discrepancy rather than attributing a cause. We report our own
protocol in full so that the comparison can be made.

\subsection{What a future collection would need}

The actionable consequence concerns label supply. A collection targeting this
task should record explicit non-stress intervals, since their absence makes
every downstream number conditional on a construction; capture onset more
finely than a retrospective interval; and prioritise electrodermal signal
quality, since four of fifteen participants here produced a flat trace on the
only responsive channel.

\section{Limitations}

The cohort is 10 of 15 female nurses at a single hospital during a pandemic,
and external validity is correspondingly narrow. Negatives are constructed, so
precision is not reported as a performance claim: when the model flags an
unlabelled window we cannot distinguish a false alarm from an unreported
episode. The class prior is not identifiable and is bracketed rather than
estimated. Our non-negative positive-unlabelled implementation uses a fixed
training budget without an inner validation split, so its negative result is
weaker evidence than the naive comparator, which shares the baseline's learner
and features and differs only in label treatment. We searched two learners and
one feature set; a sequence model over raw windows is untested. Thirty-three
rated episodes have no sensor coverage and are absent from every metric. Two
normalisation constants, the denominator floor and the output clip, are choices
we record and vary rather than defend. Ten folds is underpowered for the paired
comparisons we report, and we state spreads descriptively rather than dividing
them by the square root of the fold count. The permutation null uses twenty
draws, so its $p$-value floor is $0.048$; the observed value exceeds every draw
but a finer resolution would require more permutations.

The pretraining result carries its own qualifications. The encoder saw every
participant, so Table \ref{tab:ssl} reports transductive upper bounds rather
than estimates for an unseen person; per-fold pretraining is the clean
alternative and costs ten times the compute. One architecture and one training
budget were tried, and a larger encoder or longer schedule is untested. A
representation that is weak and one that is genuinely free of participant
identity produce the same low probe accuracy, and we separate them only by the
downstream sweep, which is indirect. The severity task is limited by 17
medium-severity episodes across the cohort, which is why we report cumulative
links and a collapsed task rather than a three-class model.

\section{Conclusion}

We intended to test whether self-supervised pretraining helps when
90.9\% of the data is unlabelled. Establishing the prerequisites
consumed the study, and the prerequisites turned out to be the result. The
labels mark a real state in one channel. A comparison group can be built that
survives falsification. A baseline built on it clears a permutation null
decisively and still recovers one episode in ten at a deployable alarm rate,
and neither a different learner nor three different treatments of the labels
move it.

Pretraining over 1,238 unlabelled hours does not move it either, and neither
does severity prediction, which needs no constructed negatives at all. We
report this as a result rather than a failure. What a future study most needs
from this one is not a score but a measurement of what was missing, and the
answer is labels: explicit non-stress intervals, finer onset resolution, and
electrodermal signal quality good enough that a participant's only responsive
channel is not flat.

\begin{ack}
We thank the authors of the source dataset for making it publicly available.
\end{ack}

\begin{thebibliography}{9}

\bibitem{elkan2008}
C.~Elkan and K.~Noto.
\newblock Learning classifiers from only positive and unlabeled data.
\newblock In \emph{Proceedings of the 14th ACM SIGKDD International Conference
  on Knowledge Discovery and Data Mining}, 2008.
\newblock DOI: 10.1145/1401890.1401920.

\bibitem{hosseini2022}
S.~Hosseini, R.~Gottumukkala, S.~Katragadda, R.~T. Bhupatiraju, Z.~Ashkar,
  C.~W. Borst, and K.~Cochran.
\newblock A multimodal sensor dataset for continuous stress detection of nurses
  in a hospital.
\newblock \emph{Scientific Data}, 9:255, 2022.
\newblock DOI: 10.1038/s41597-022-01361-y.

\bibitem{kiryo2017}
R.~Kiryo, G.~Niu, M.~C. du~Plessis, and M.~Sugiyama.
\newblock Positive-unlabeled learning with non-negative risk estimator.
\newblock In \emph{Advances in Neural Information Processing Systems 30}, 2017.
\newblock arXiv:1703.00593.

\bibitem{mathur2024}
A.~Mathur et~al.
\newblock Body sensor-based multimodal nurse stress detection using machine
  learning.
\newblock In \emph{2024 16th International Conference on Communication Systems
  and Networks (COMSNETS)}, 2024.

\bibitem{munoz2015}
M.~L. Munoz, A.~van~Roon, H.~Riese, C.~Thio, E.~Oostenbroek, I.~Westrik,
  E.~J.~C. de~Geus, R.~Gansevoort, J.~Lefrandt, I.~M. Nolte, and H.~Snieder.
\newblock Validity of (ultra-)short recordings for heart rate variability
  measurements.
\newblock \emph{PLOS ONE}, 10(9):e0138921, 2015.
\newblock DOI: 10.1371/journal.pone.0138921.

\bibitem{orini2023}
M.~Orini et~al.
\newblock Long-term association of ultra-short heart rate variability with
  cardiovascular events.
\newblock \emph{Scientific Reports}, 13, 2023.

\end{thebibliography}


\newpage

\appendix

\section{Broader impacts}

The application is monitoring of employees during work, and that is worth
naming plainly. A reliable stress detector deployed in a hospital could direct
support toward staff who need it, which is the motivating case. The same
instrument could as easily support surveillance, scheduling penalties, or
insurance decisions, and the people measured are in a weak position to refuse.
Consent to be monitored by an employer is not freely given in the way consent
to a research study is.

Our result bears on this directly. A detector recovering one episode in ten at
a deployable alarm rate is not fit for any of those uses, including the
benevolent one, and two of ten participants are undetectable at any usable
threshold. We would regard deployment on the strength of results at this level
as harmful, and we note that a published figure of $F_1 = 0.99$ on this same
dataset could be read as licensing exactly that. Reporting the achievable
number honestly is the contribution most relevant to impact here.

The cohort is ten female nurses at one hospital during a pandemic. Any
performance claim generalised beyond that population would be unsupported, and
differential performance across groups cannot be assessed at this sample size.

\section{Compute} All experiments run on a single laptop CPU with no GPU.
Feature extraction over 387 sessions takes 107 s; the complete set
of reported experiments completes in about half an hour beyond feature
extraction. Pretraining over 68k windows takes 192 s for twelve epochs, and the
label-efficiency sweep, which is 120 model fits, takes roughly thirteen
minutes. The audit that precedes
them reads the 3.5 GB archive once per section and is likewise
CPU-bound. No experiment reported here was preceded by a larger unreported
sweep.

\section{Per-participant cohort and negative construction}
\label{app:cohort}

The modelling set behind every reported result: the ten retained
participants at the fixed negative-to-positive ratio of 1.31 adopted in
Section 4.3, for negative-sampling seed 0. A ratio below the 1.31 target
means the eligible negative pool was exhausted before the target was met;
the lowest, at 1.03, sets the residual balance spread of
$1.28\times$ quoted in Table 3.

\begin{table}[h]
\centering
\begin{tabular}{lrrr}
\toprule
Participant & Positive windows & Negative windows & Ratio \\
\midrule
8B & 74 & 76 & 1.03 \\
5C & 68 & 84 & 1.24 \\
94 & 34 & 44 & 1.29 \\
7A & 248 & 324 & 1.31 \\
6B & 182 & 238 & 1.31 \\
E4 & 310 & 406 & 1.31 \\
F5 & 116 & 152 & 1.31 \\
83 & 314 & 412 & 1.31 \\
15 & 79 & 104 & 1.32 \\
BG & 154 & 203 & 1.32 \\
\midrule
Total & 1579 & 2043 & 1.29 \\
\bottomrule
\end{tabular}
\caption{Cohort and negative construction at the adopted ratio of 1.31, sorted by achieved ratio.}
\end{table}

\section{Exclusion cascade}
\label{app:attrition}

Window counts through the exclusion rules, in the order fixed in advance.
Rule numbering follows the analysis plan; rules that remove no windows at
the window level are omitted.

\begin{table}[h]
\centering
\begin{tabular}{cp{5.8cm}rrr}
\toprule
\# & Rule & Entering & Removed & Surviving \\
\midrule
1 & sessions $<$ 5 min, before non-wear removal & 37252 & 32 & 37220 \\
2 & non-wear windows & 37220 & 87 & 37133 \\
3 & dead-EDA sessions & 37133 & 598 & 36535 \\
8 & participant 6D, no eligible negatives & 36535 & 681 & 35854 \\
9 & flat-EDA participants (four) & 35854 & 11520 & 24334 \\
\bottomrule
\end{tabular}
\caption{Exclusion cascade at the window level.}
\end{table}

\section{Per-fold results}
\label{app:perfold}

Leave-one-participant-out folds, averaged over three negative-sampling
seeds, with the across-seed range. Episode recall is at one false alarm
per worn hour. These are the folds behind the aggregate figures in
Section 5; we report them individually because the spread is wide and the
two lowest-count folds are not comparable to the rest.

\begin{table}[h]
\centering
\begin{tabular}{lrrcr}
\toprule
Participant & Episodes & AUC & AUC range & Episode recall \\
\midrule
5C & 4 & 0.388 & 0.365--0.401 & 0.000 \\
BG & 13 & 0.633 & 0.622--0.646 & 0.000 \\
6B & 10 & 0.682 & 0.671--0.701 & 0.200 \\
15 & 10 & 0.811 & 0.794--0.828 & 0.067 \\
E4 & 23 & 0.813 & 0.804--0.825 & 0.232 \\
94 & 5 & 0.835 & 0.816--0.865 & 0.133 \\
83 & 15 & 0.841 & 0.828--0.861 & 0.089 \\
7A & 23 & 0.847 & 0.827--0.867 & 0.029 \\
F5 & 23 & 0.882 & 0.867--0.898 & 0.159 \\
8B & 12 & 0.885 & 0.858--0.909 & 0.028 \\
\midrule
Mean & 138 & 0.762 & & 0.094 \\
\bottomrule
\end{tabular}
\caption{Per-fold AUC and episode recall, seed-averaged. The fold range quoted in the main text, 0.388 to 0.885, is the range of this column.}
\end{table}

\section{Operating points}
\label{app:operating}

The alarm-rate trade-off, averaged over three seeds. The paper reports the
one-alarm-per-worn-hour row; the others are given so the shape of the
curve is visible rather than asserted.

\begin{table}[h]
\centering
\begin{tabular}{rrrrc}
\toprule
False alarms/h & FPR & Window recall & Episode recall & Range \\
\midrule
0.5 & 0.0176 & 0.087 & 0.029 & 0.029--0.029 \\
1.0 & 0.0357 & 0.155 & 0.106 & 0.058--0.130 \\
2.0 & 0.0715 & 0.292 & 0.280 & 0.261--0.319 \\
\bottomrule
\end{tabular}
\caption{Operating points over 138 episodes with usable coverage.}
\end{table}


\newpage
\section*{NeurIPS Paper Checklist}

\begin{enumerate}

\item {\bf Claims}
    \item[] Question: Do the main claims made in the abstract and introduction accurately reflect the paper's contributions and scope?
    \item[] Answer: \answerYes{}
    \item[] Justification: The abstract states the negative result and its boundary explicitly: we eliminated the label-treatment axis and not the representation axis, and Section 7.1 repeats this rather than generalising past it. The intended contribution, an account of the prerequisites for self-supervision on sparsely labelled physiological data, matches what is delivered. Every quantity quoted in the abstract appears with its seed range where one exists.

\item {\bf Limitations}
    \item[] Question: Does the paper discuss the limitations of the work performed by the authors?
    \item[] Answer: \answerYes{}
    \item[] Justification: Section 8 is a dedicated limitations section. It names the narrow cohort, the constructed rather than observed negative class and the consequent inestimability of precision, the unidentifiable class prior, the undertuned positive-unlabelled arm and why its negative result is weaker evidence than the naive comparator, the two-learner search space, the 33 episodes with no sensor coverage, the two normalisation constants chosen rather than derived, the underpowered ten-fold comparisons, and the resolution floor of a twenty-draw permutation test.

\item {\bf Theory assumptions and proofs}
    \item[] Question: For each theoretical result, does the paper provide the full set of assumptions and a complete (and correct) proof?
    \item[] Answer: \answerNA{}
    \item[] Justification: The paper contains no theoretical results. The one formal claim it relies on, that the class prior is not identifiable from positive-unlabelled data, is cited rather than proved.

    \item {\bf Experimental result reproducibility}
    \item[] Question: Does the paper fully disclose all the information needed to reproduce the main experimental results of the paper to the extent that it affects the main claims and/or conclusions of the paper (regardless of whether the code and data are provided or not)?
    \item[] Answer: \answerYes{}
    \item[] Justification: The source dataset is public and cited. The paper specifies the window length and hop, the label rule and its threshold, the exclusion rules and the order in which they are applied, the four eligibility conditions for negatives, the matching procedure and its binning, the negative ratio and how it was chosen, the number of sampling seeds, the feature set by channel, the normalisation and its two constants, the learner, and the fold structure. Where a choice could have gone otherwise, the alternative is reported alongside it rather than omitted.

\item {\bf Open access to data and code}
    \item[] Question: Does the paper provide open access to the data and code, with sufficient instructions to faithfully reproduce the main experimental results, as described in supplemental material?
    \item[] Answer: \answerNo{}
    \item[] Justification: The dataset is public and cited by DOI. The analysis code is not released with this submission, because the repository would identify the authors and the review is double-blind. The protocol is specified in full in Sections 2 through 6 with the intent that it be reimplementable from the paper alone. Code will be released on acceptance or on publication elsewhere.

\item {\bf Experimental setting/details}
    \item[] Question: Does the paper specify all the training and test details (e.g., data splits, hyperparameters, how they were chosen, type of optimizer) necessary to understand the results?
    \item[] Answer: \answerYes{}
    \item[] Justification: Sections 2, 4 and 5 give the splits (leave-one-participant-out, grouped on participant, with the reason stated), the learner and its class weighting, the normalisation window and its floor and clip, the negative ratio and the sensitivity analysis over it, and the class prior grid for the positive-unlabelled arm. Choices fixed before results were seen are identified as such, including the episode-detection fraction and the decision rule in Section 6.

\item {\bf Experiment statistical significance}
    \item[] Question: Does the paper report error bars suitably and correctly defined or other appropriate information about the statistical significance of the experiments?
    \item[] Answer: \answerYes{}
    \item[] Justification: Every quantity computed on the constructed negatives is reported with its range across three negative-sampling seeds, and that range is what the variability captures: the stochastic draw of the comparison group, holding folds, features and learner fixed. The onset analysis reports a sign test and a Wilcoxon signed-rank test rather than a parametric interval, because the underlying distribution is heavy-tailed for reasons the paper explains. The baseline is compared against a twenty-draw within-participant permutation null, reported as mean, standard deviation and maximum. Where the seed range is comparable to the estimate itself, as it is for episode recall, the paper says so at every occurrence rather than only once. Fold-level results are reported individually and never averaged, since two of ten folds behave qualitatively differently from the rest.

\item {\bf Experiments compute resources}
    \item[] Question: For each experiment, does the paper provide sufficient information on the computer resources (type of compute workers, memory, time of execution) needed to reproduce the experiments?
    \item[] Answer: \answerYes{}
    \item[] Justification: A compute section in Appendix B states that all experiments run on a single laptop CPU with no GPU, gives the feature-extraction time over 387 sessions, and states that the complete set of reported experiments finishes in under an hour. It also records that no reported experiment was preceded by a larger unreported sweep.

\item {\bf Code of ethics}
    \item[] Question: Does the research conducted in the paper conform, in every respect, with the NeurIPS Code of Ethics \url{https://neurips.cc/public/EthicsGuidelines}?
    \item[] Answer: \answerYes{}
    \item[] Justification: The work is a secondary analysis of a public, de-identified dataset whose original collection was ethically reviewed by the depositing institution. No new data were collected and no participant was contacted. Anonymity is preserved in this submission.

\item {\bf Broader impacts}
    \item[] Question: Does the paper discuss both potential positive societal impacts and negative societal impacts of the work performed?
    \item[] Answer: \answerYes{}
    \item[] Justification: Appendix A addresses this. The positive case is directing support toward staff who need it. The negative case is that the same instrument supports surveillance, scheduling penalties, or insurance decisions, and that employees are poorly placed to refuse monitoring by an employer. The section argues that a detector performing at the level reported here is unfit for any of these uses including the benevolent one, and notes that a published figure of $F_1 = 0.99$ on this dataset could be read as licensing deployment that our results do not support.

\item {\bf Safeguards}
    \item[] Question: Does the paper describe safeguards that have been put in place for responsible release of data or models that have a high risk for misuse (e.g., pre-trained language models, image generators, or scraped datasets)?
    \item[] Answer: \answerNA{}
    \item[] Justification: No model, dataset, or scraped resource is released. The trained models are not distributed and, on the paper's own argument, are not fit for deployment.

\item {\bf Licenses for existing assets}
    \item[] Question: Are the creators or original owners of assets (e.g., code, data, models), used in the paper, properly credited and are the license and terms of use explicitly mentioned and properly respected?
    \item[] Answer: \answerYes{}
    \item[] Justification: The dataset is credited to its depositors and cited both as a Dryad deposit by DOI and as the accompanying data descriptor. The deposit landing page states that content is licensed for reuse but does not display a specific license designation, so we describe it as publicly available under the repository's terms rather than naming a license we cannot verify. Software used is open source: scikit-learn, NumPy, pandas, and NeuroKit2 for electrodermal response detection.

\item {\bf New assets}
    \item[] Question: Are new assets introduced in the paper well documented and is the documentation provided alongside the assets?
    \item[] Answer: \answerNA{}
    \item[] Justification: No new assets are released with this submission.

\item {\bf Crowdsourcing and research with human subjects}
    \item[] Question: For crowdsourcing experiments and research with human subjects, does the paper include the full text of instructions given to participants and screenshots, if applicable, as well as details about compensation (if any)?
    \item[] Answer: \answerNA{}
    \item[] Justification: This is a secondary analysis of an existing public dataset. We did not recruit, instruct, compensate, or interact with any participant. Instructions given to the original participants and the survey instrument are documented by the depositors in the cited data descriptor.

\item {\bf Institutional review board (IRB) approvals or equivalent for research with human subjects}
    \item[] Question: Does the paper describe potential risks incurred by study participants, whether such risks were disclosed to the subjects, and whether Institutional Review Board (IRB) approvals (or an equivalent approval/review based on the requirements of your country or institution) were obtained?
    \item[] Answer: \answerNA{}
    \item[] Justification: No human subjects research was conducted for this paper. The original collection was reviewed and approved by the depositing institution's review board, as reported in the cited data descriptor; we rely on that record and did not seek separate approval for secondary analysis of de-identified public data.

\item {\bf Declaration of LLM usage}
    \item[] Question: Does the paper describe the usage of LLMs if it is an important, original, or non-standard component of the core methods in this research? Note that if the LLM is used only for writing, editing, or formatting purposes and does \emph{not} impact the core methodology, scientific rigor, or originality of the research, declaration is not required.
    \item[] Answer: \answerNA{}
    \item[] Justification: No large language model is a component of the method. The pipeline is signal processing, feature extraction, and gradient boosting or logistic regression; no LLM runs at training or inference time, and no reported result depends on one. Separately, and as voluntary disclosure rather than because the question requires it: a coding assistant was used while implementing the analysis scripts and drafting the manuscript. Every reported number was produced by those scripts and re-derived from the frozen result tables, and every reference was checked against Crossref, arXiv, or a literature index rather than generated.

\end{enumerate}


\end{document}
